\documentclass[12pt]{article}
\usepackage[utf8]{inputenc}
\usepackage{geometry}
\geometry{a4paper,scale=0.75}
\linespread{1.5}
\usepackage{graphicx} 
\usepackage{float} 
\usepackage{subcaption} 
\usepackage{enumerate}
\usepackage{amsmath}
\usepackage{array}
\newcolumntype{P}[1]{>{\raggedright\arraybackslash}p{#1}}
\usepackage{booktabs}
\usepackage{multirow}
\usepackage{amsfonts}
\usepackage[english]{babel}
\usepackage{amsthm}
\usepackage{dcolumn}
\usepackage{multicol}
\usepackage{stfloats}
\usepackage{placeins}
\usepackage{lscape}
\usepackage[figuresright]{rotating}
\RequirePackage{pdflscape}
\usepackage[toc,page]{appendix}
\usepackage{geometry}
\usepackage{longtable}
\usepackage{comment}
\usepackage{xcolor}
% \usepackage{amsmath}
\usepackage{amssymb}
\usepackage{empheq}
\usepackage{newtxtext,newtxmath}

% -------- enumerated sub-labels (a), (b), … --
\usepackage{enumitem}
\setlist[enumerate,1]{label=(\alph*),ref=\alph*}
% ---------------------------------------------

\usepackage{hyperref}
\hypersetup{hidelinks,
	colorlinks=true,
	allcolors=black,
	pdfstartview=Fit,
	breaklinks=true}
\usepackage{csquotes}
\usepackage{natbib}
\newtheorem{definition}{Definition}
\newtheorem{theorem}{Theorem}
%\newtheorem{proposition}[theorem]{Proposition}
\newtheorem{proposition}{Proposition}[section]
%\newtheorem{lemma}[theorem]{Lemma}
\newtheorem{lemma}{Lemma}[section]
%\newtheorem{corollary}[theorem]{Corollary}
\newtheorem{corollary}{Corollary}[section]
\newtheorem*{remark}{Remark}
\newtheorem{example}{Example}
\newtheorem{exercise}{Exercise}[section]
\numberwithin{equation}{section}
\newtheorem{assumption}{Assumption}[section] % number within sections

\title{Production Networks and Formation of Aggregate Uncertainty}
\author{
Harlly Zhou\footnotemark[1]\\\texttt{hzhouci@ust.hk}
}
\date{\today}

\begin{document}

\begin{titlepage}
    \begin{center}
        {\LARGE\bfseries Production Networks and Formation of Aggregate Uncertainty\par}
        \vspace{1.5em}
        {\large
        Harlly Zhou\footnotemark[1] (\texttt{hzhouci@ust.hk})\\
        }
        \vspace{1em}
        \today
    \end{center}
    \vspace{2em}
    \begin{abstract}
        \sloppy
        \noindent 
         \vspace{1em}

        \noindent\textbf{JEL Codes:} 
        \noindent\textbf{Keywords:} 
        \fussy
    \end{abstract}
    \thispagestyle{empty}
    \setcounter{page}{0}
\end{titlepage}


\section{Introduction}\label{sec:introduction}


\citep{acemogluEndogenousProductionNetworks2020}

\section{Benchmark Model}\label{sec:benchmark}
In the benchmark model, we assume that all the agents have full information with rational expectations (FIRE). The friction in this economy is the financial friction in the form of borrowing constraint.

Time is discrete and indexed by $t=0, 1, 2, \cdots$. There are representative agents with unit mass who consume a consumption aggregator $C_t$, make nominal deposits $D_t$ and supplies labour $N_t$, a representative firm in each sector $i\in\{1, 2, \cdots, N\}$ \footnote{In the following, we will use the sector index $i$ and firm $i$ interchangeably.} that produces a differentiatiated good $y_{i,t}$ using labour $n_{i,t}$ and intermediate input $x_{ij, t}$ from sectors $j$, competitive financial intermediaries that takes deposits from households and supplies credit to firms, and a monetary authority that sets the nominal interest rate $i_t$.

Interest rates are dated by when they are set, but need not apply over the same window. The policy rate $i_t$ is set at the beginning of period $t$ and remunerates deposits from the beginning of period $t+1$ to its end (so deposits $D_t$ earn $i_t$ in period $t+1$). The intra-period loan rate $i^L_{i,t}$ is set at the beginning of period $t$ and applies from the beginning of period $t$ to its end; it is priced off the predetermined funding rate $i_{t-1}$ that prevails over that same window.

The timeline within a period $t$ is as follows. At the beginning of the period, the monetary authority sets the nominal interest rate $i_t$. The financial intermediary charges firm $i$ the loan rate $i^L_{i,t}$, and firms borrow, hire labour, and buy intermediate inputs. At the end of the period, productivity realizes, firms produce and sell goods and repay $\left(1+i^L_{i,t}\right)b_{i,t}$, and households consume, collect the maturity payoff $\left(1+i_{t-1}\right)$ on deposits $D_{t-1}$, and choose new deposits $D_t$.

\subsection{Environment}
\subsubsection{Households}

Households gain utility from a consumption aggregator $C_t$ and disutility from labour supply $N_t$:
\begin{align}\label{eq:hh_utility}
    U(C_t, N_t) = \log C_t - \frac{\chi}{1+\varphi} N_t^{1+\varphi},
\end{align}
with $\chi > 0$ and $\varphi>0$, where the consumption aggregator is given by a CES aggregator of the differentiated goods $c_{i,t}$:
\begin{align}\label{eq:cons_aggregator}
    C_t = \left(\sum_{i=1}^N \rho_i^{\frac{1}{\xi}} c_{i,t}^{\frac{\xi-1}{\xi}}\right)^{\frac{\xi}{\xi-1}}
\end{align}
with $\rho_i \in (0,1)$ the demand weight of goods $i$ satisfying $\sum_{i=1}^N \rho_i = 1$ and $\xi>0$ the elasticity of substitution (EoS) between goods. As a result, the price level is given by
\begin{align}\label{eq:price_level}
    P_t = \left(\sum_{i=1}^N \rho_i p_{i,t}^{1-\xi}\right)^{\frac{1}{1-\xi}}.
\end{align}
where $p_{i,t}$ is the price of good $i$.

They are subject to a budget constraint:
\begin{align}\label{eq:hh_budget}
    P_t C_t + D_t \leq W_t N_t + \left(1+i^D_{t-1}\right) D_{t-1} + \Pi^y_t + \Pi^f_t,
\end{align}
where $W_t$ is the nominal wage, $i^D_{t-1}$ is the nominal interest rate on deposits made in period $t-1$ (paid in period $t$), $\Pi^y_t$ is the total nominal profit of all the firms, and $\Pi^f_t$ is the total nominal profit of the financial intermediary which is allowed to be negative to reflect possible defaults.

\subsubsection{Firms}
Firm $i$ uses intermediate inputs from other sectors so that the economy has a network, denoted by an $N \times N$ matrix $\Omega_t = \left(\omega_{ij,t}\right)_{i,j=1}^N$ where the $(i,j)$-th entry $\omega_{ij,t}$ denotes the expenditure share of firm $i$ on input $j$. The corresponding quantity of $j$ used to produce $i$ is $x_{ij,t}$. Firms take input prices as given. For simplicity, assume that prices are pre-determined before productivity $A_{i,t}$ realizes.

The firm produces its output $y_{i,t}$ using the following production function:
\begin{align}\label{eq:firm_prod_fcn}
    y_{i,t} = A_{i,t} \underbrace{n_{i,t}^{\alpha_{i,t}} \prod_{j=1}^N x_{ij, t}^{\omega_{ij,t}}}_{\equiv \tilde{y}_{i,t}}
\end{align}
with $\alpha_{i,t} + \sum_{j=1}^N \omega_{ij, t} = 1$, where $A_{i,t}$ is the stochastic productivity
\begin{align}\label{eq:firm_productivity}
    \log A_{i,t} = \bar{a}_i + \sigma_{i,t-1} \epsilon_{i,t}.
\end{align}

We look at the components of productivity in \eqref{eq:firm_productivity} one by one. $\bar{a}_i$ denotes the sector-specific average productivity. $\sigma_{i,t}$ denotes the standard deviation scale of log productivity. $\epsilon_{i,t}\sim\mathcal{N}(0,1)$ with joint normal distribution ${\bf \epsilon}_t = (\epsilon_{1,t}, \epsilon_{2,t}, \cdots, \epsilon_{N,t})' \sim \mathcal{N}\left({\bf 0}, \Sigma_{\epsilon,t}\right)$ denotes the number of units of standard deviation. 

Assume that uncertainty shocks are independent, \text{i.e.,} $\sigma_{i,t}$'s are independent among sectors and are independent of level shocks $\epsilon_{i,t}$. Assume that the uncertainty shocks follow: 
\begin{align}\label{eq:uncertainty_shock_process}
    \log \sigma^2_{i,t} = (1-\phi)\log \bar{\sigma}^2_i + \phi \log \sigma^2_{i,t-1} + v_{i,t},
\end{align}
where $\phi\in(0,1)$ and $v_{i,t} \sim \mathcal{N}(0,\sigma^2_v)$ is the uncertainty shock. This decomposition isolates the uncertainty shock from the level shock: a shock to $\epsilon_{i,t}$ changes the realized productivity but not the distribution, while a shock to $\sigma_{i,t}$ changes the future conditional distribution of productivity. To be more precise, since $\sigma_{i,t} \perp \epsilon_{i,t}$, 
\begin{align*}
    \operatorname{Var}_{t-1}\left(\log A_{i,t}\right) = \operatorname{Var}_{t-1}\left(\sigma_{i,t-1} \epsilon_{i,t}\right) = \sigma_{i,t-1}^2.
\end{align*}


The following lemma provides the closed-form solution to the conditional mean and standard deviation of productivity.
\begin{lemma}[Conditional Mean and Standard Deviation of Productivity]\label{lemma:conditional_mean_std_productivity_benchmark}
    Given the information at the end of period $t-1$, productivity at period $t$ is log-normally distributed with the conditional mean and standard deviation given by
    \begin{align*}
        \mu^A_{i,t|t-1} = e^{\bar{a}_i + \frac{1}{2}\sigma_{i,t-1}^2},
        \quad
        \sigma^A_{i,t|t-1} = \mu^A_{i,t|t-1}\sqrt{e^{\sigma_{i,t-1}^2} - 1}.
    \end{align*}
\end{lemma}
\begin{proof}
    Conditional on the information at the end of period $t-1$, \eqref{eq:uncertainty_shock_process} implies that the log of productivity has mean $\bar{a}_i$ and standard deviation $\sigma_{i,t-1}$. The remaining follows from the properties of log-normal distribution.
\end{proof}

At the beginning of each period, each firm $i$ faces labour costs $W_t n_{i,t}$ and intermediate input costs $\sum_{j=1}^N p_{j,t}x_{ij, t}$. Therefore, it needs to borrow $b_{i,t}$ from the financial intermediary to cover the costs:
\begin{align}\label{eq:firm_financial_constraint}
    W_t n_{i,t} + \sum_{j=1}^N p_{j,t}x_{ij, t} \leq  b_{i,t},
\end{align}
subject to a borrowing constraint, which will be discussed soon in the Section \ref{sec:financial_intermediary}. At the end of the period, the firm has a revenue $R_{i,t} = p_{i,t} y_{i,t}$ and repays the debt $b_{i,t}$ together with interest at the intra-period loan rate $i^L_{i,t}$. Its nominal profit 
\begin{align*}
\Pi^y_{i,t} = \max\left\{R_{i,t} - \left(1 + i^L_{i,t}\right) b_{i,t}, 0\right\}
\end{align*}
will be distributed to the households.


\subsubsection{Financial Intermediary}\label{sec:financial_intermediary}

Competitive financial intermediaries take the deposit rate as tied to the policy rate $i^D_t=i_t$, so deposits $D_t$ earn $i_t$ in period $t+1$. Period-$t$ working-capital loans are funded by outstanding nominal deposits $D_{t-1}$ and are priced off the predetermined funding rate $i_{t-1}$ using a common risk-pricing rule:
\begin{align*}
    i^L_{t}
    =
    i_{t-1}
    +
    s_{i,t},
\end{align*}
where $s_{i,t}$ is the credit spread satisfying
\begin{align}\label{eq:credit_spread_rule}
    s_{i,t}
    =
    \bar{s}
    +
    \psi U_t,
\end{align}
for some $\psi>0$, where $U_t$ is the aggregate uncertainty measure as an weighted average of the firm-level uncertainty measures à la \citep{juradoMeasuringUncertainty2015} and \citep{ludvigsonUncertaintyBusinessCycles2021} defined as the variance \footnote{The original papers defined the uncertainty measure in terms of standard deviation. We use variance here for consistency with the assumed process of the uncertainty shock.} of the forecast error of an objective variable, namely the logarithm of productivity in this paper:
\begin{align*}
    U_t := \sum_{i=1}^N \omega^U_{i,t-1} \sigma^2_{i,t-1},
\end{align*}
where the weight is defined as the share of firm $i$'s nominal output in the aggregate output:
\begin{align*}
    \omega^U_{i,t-1} = \frac{p_{i,t} y_{i,t}}{\sum_{i=1}^N p_{i,t} y_{i,t}}.
\end{align*}

The intermediary lends $\ell_{i,t}$ to firm $i$ with the nominal balance-sheet constraint
\begin{align}\label{eq:balance_sheet_constraint}
    \sum_{i=1}^N \ell_{i,t}
    \leq
    D_{t-1}.
\end{align}
Its nominal profit $\Pi^f_t$ is given by
\begin{align*}
    \Pi^f_t = \sum_{i=1}^N \left[\min\left\{R_{i,t}, (1+i^L_{t})\ell_{i,t}\right\} - \ell_{i,t} \right]  - i_{t-1} D_{t-1}.
\end{align*}
For each firm $i$, the intermediary imposes the borrowing constraint 
\begin{align}
    \ell_{i,t}
    \leq
    \kappa \mu^R_{i,t|t-1} \exp\left(-\theta \sigma^2_{i,t-1}\right),
    \label{eq:borrowing_constraint}
\end{align}
where
\begin{align*}
    \mu^R_{i,t|t-1}
    =
    \mathbb{E}_{t-1}\left[R_{i,t}\right] = p_{i,t} \tilde{y}_{i,t} \mu^A_{i,t|t-1},
\end{align*}
is the mean of revenue conditional on information available at the end of period $t-1$. The parameter $\kappa>0$ captures the pledgeability of expected revenue, and $\theta>0$ captures the sensitivity of credit capacity to revenue risk. This constraint prevents the intermediary from severe credit risk, but it cannot eliminate default. Moreover, the realized borrowing amount will not exceed the lending amount:
\begin{align*}
    b_{i,t} \leq \ell_{i,t}, \forall i.
\end{align*}


We will not explicitly solve for the financial intermediary's optimization problem in the benchmark model. We show that the credit-spread rule and the borrowing constraint can be interpreted as a reduced-form representation of competitive risk-based lending.

The financial intermediary's optimization problem is to maximize its nominal profit subject to a default-tolerance rule:
\begin{align*}
    \max_{b_{i,t}} \mathbb{E}_{t-1}\Pi^f_t,
\end{align*}
subject to
\begin{align}\label{eq:default_tolerance_rule}
    \mathbb{P}_{t-1} \left( R_{i,t} < (1+i^L_{t})\ell_{i,t} \right) \leq \delta,
\end{align}
for some $\delta \in (0,1)$, and the balance sheet constraint \eqref{eq:balance_sheet_constraint}.

We first derive the implication of the default-tolerance rule. Note that $R_{i,t} = p_{i,t} \tilde{y}_{i,t} A_{i,t}$. Therefore, the rule implies that
\begin{align*}
    \mathbb{P}_{t-1} \left( R_{i,t} < (1+i^L_{t})\ell_{i,t} \right) &= \mathbb{P}_{t-1} \left( \log A_{i,t} < \log\left((1+i^L_{t})\ell_{i,t}\right) - \log p_{i,t}\tilde{y}_{i,t} \right) \\
    &= 1-\Phi(q_{1,i})\\
    &\leq \delta,
\end{align*}
where
\begin{align*}
    q_{1,i} = \frac{\bar{a}_i - \log\left((1+i^L_{t})\ell_{i,t}\right) + \log \left(p_{i,t}\tilde{y}_{i,t}\right)}{\sigma_{i,t-1}}.
\end{align*}
The exact default-tolerance rule yields
\begin{align*}
    \ell_{i,t}
    \leq
    \frac{\mu^R_{i,t|t-1}}{1+i^L_{i,t}}
    \exp\left(
    -\frac{1}{2}\sigma^2_{i,t-1}
    -
    z_{1-\delta}\sigma_{i,t-1}
    \right).
\end{align*}
This motivates the reduced-form borrowing-capacity constraint \eqref{eq:borrowing_constraint} where $\kappa>0$ captures pledgeability and average funding-cost effects, while $\theta>0$ captures the sensitivity of credit capacity to revenue risk.

Now consider the objective function:
\begin{align}\label{eq:exp_fina_intermediary_profit}
    \mathbb{E}_{t-1}\Pi^f_t = \sum_{i=1}^N \left[p_{i,t}\mu^A_{i,t|t-1}\tilde{y}_{i,t}\left(1-\Phi(q_{2,i})\right) + (1+i^L_{i,t})\ell_{i,t} \Phi(q_{1,i}) - \ell_{i,t}\right] - i_{t-1}D_{t-1},
\end{align}
where
\begin{align*}
    q_{2,i} &= \frac{\bar{a}_i + \sigma^2_{i,t-1} - \log\left((1+i^L_{t})\ell_{i,t}\right) + \log \left(p_{i,t}\tilde{y}_{i,t}\right)}{\sigma_{i,t-1}} = q_{1,i} + \sigma_{i,t-1}.
\end{align*}

Here we pause a bit for a technical lemma that is directly obtained from the definitions of $q_{1,i}$ and $q_{2,i}$.
\begin{lemma}[Densities at $q_{1,i}$ and $q_{2,i}$]\label{lemma:pdf_q1_q2}
    The probability densities of $q_{1,i}$ and $q_{2,i}$ have the following relationship:
    \begin{align*}
        \frac{\phi(q_{2,i})}{\phi(q_{1,i})} = \frac{\left(1+i^L_{t}\right)\ell_{i,t}}{p_{i,t}\mu^A_{i,t|t-1}\tilde{y}_{i,t}}.
    \end{align*}
\end{lemma}
\begin{proof}
    Note that $q_{2,i} = q_{1,i} + \sigma_{i,t-1}$. Therefore,
    \begin{align*}
        \frac{\phi(q_{2,i})}{\phi(q_{1,i})} = \exp \left(-q_{1,i}\sigma_{i,t-1} -\frac{1}{2} \sigma_{i,t-1}^2 \right).
    \end{align*}
    Substituting the expression of $q_{1,i}$ into the expression yields
    \begin{align*}
        \frac{\phi(q_{2,i})}{\phi(q_{1,i})} &= \exp \left(-\bar{a}_i + \log\left((1+i^L_{t})\ell_{i,t}\right) - \log \left(p_{i,t}\tilde{y}_{i,t}\right) - \frac{1}{2}\sigma_{i,t-1}^2 \right)\\
        &= \frac{\left(1+i^L_{t}\right)\ell_{i,t}}{p_{i,t}\tilde{y}_{i,t}} \exp \left(-\bar{a}_i - \frac{1}{2}\sigma_{i,t-1}^2 \right).
    \end{align*}
    By Lemma \ref{lemma:conditional_mean_std_productivity_benchmark}, $\exp \left(-\bar{a}_i - \frac{1}{2}\sigma_{i,t-1}^2 \right) = \left(\mu^A_{i, t|t-1}\right)^{-1}$. We are done. 
\end{proof}

\begin{corollary}[Default-tolerance constraint binds]\label{corollary:default_tolerance_constraint_binds}
    When the default-tolerance constraint binds, $q_{1,i} = z_{1-\delta}$ for all $i$.
\end{corollary}
\begin{proof}
    Suppose that the default-tolerance constriant binds. Then we have
    \begin{align}\label{eq:default_tolerance_constraint_binds}
        (1+i^L_{t})\ell_{i,t} = p_{i,t}\mu^A_{i,t|t-1}\tilde{y}_{i,t}\exp\left(-\frac{1}{2}\sigma^2_{i,t-1} - z_{1-\delta}\sigma_{i,t-1}\right).
    \end{align}
    By Lemma \ref{lemma:pdf_q1_q2}, we have
    \begin{align*}
        \exp\left(-\frac{1}{2}\sigma^2_{i,t-1} - z_{1-\delta}\sigma_{i,t-1}\right) = \frac{\phi(q_{2,i})}{\phi(q_{1,i})} = \exp\left(-\frac{1}{2}\sigma^2_{i,t-1} - q_{1,i}\sigma_{i,t-1}\right).
    \end{align*}
    It follows from the the monotonicty of the exponential functions that $z_{1-\delta} = q_{1,i}$.
\end{proof}

The derivation below shows that competitive risk-based lending implies a loan rate that is increasing in revenue risk. We use the linear spread rule \eqref{eq:credit_spread_rule} as a reduced-form approximation to this structural pricing relation.

Since financial intermediaries are perfectly competitive, they should have zero expected profit for each loan $i$ made. The potential realized loss will be borne by household consumption. Suppose that the default-tolerance constraint binds. 
Thien equation \eqref{eq:default_tolerance_constraint_binds} and the zero-expected profit condition for loan $i$ together imply that
\begin{align*}
    p_{i,t}\mu^A_{i,t|t-1}\tilde{y}_{i,t}\left(1-\Phi\left(z_{1-\delta}+\sigma_{i,t-1}\right)\right) + (1+i^L_{t})\ell_{i,t}(1-\delta) = (1+i_{t-1})\ell_{i,t}.
\end{align*}
Further simplifying the expression yields
\begin{align*}
    1 + i^L_{t} = (1+i_{t-1})\left[\exp\left(\frac{1}{2}\sigma^2_{i,t-1}+z_{1-\delta}\sigma_{i,t-1}\right)\left(1-\Phi\left(z_{1-\delta}+\sigma_{i,t-1}\right)\right) + (1-\delta)\right]^{-1}
\end{align*}
We now show that the multiplier to $1+i_{t-1}$ is greater than or equal to 1. Note that
\begin{align*}
    &\exp\left(\frac{1}{2}\sigma^2_{i,t-1}+z_{1-\delta}\sigma_{i,t-1}\right)\left(1-\Phi\left(z_{1-\delta}+\sigma_{i,t-1}\right)\right)\\
    = \,\,&\exp\left(\frac{1}{2}\sigma^2_{i,t-1}+z_{1-\delta}\sigma_{i,t-1}\right) \int^{+\infty}_{z_{1-\delta}+\sigma_{i,t-1}}\phi(u)du\\
    = \,\,&\exp\left(\frac{1}{2}\sigma^2_{i,t-1}+z_{1-\delta}\sigma_{i,t-1}\right) \int^{+\infty}_{z_{1-\delta}}\phi(x+\sigma_{i,t-1})dx\\
    = \,\,&\int^{+\infty}_{z_{1-\delta}}\phi(x) \exp\left(-\sigma(x-z_{1-\delta})\right)dx\\
    \leq \,\,&\int^{+\infty}_{z_{1-\delta}}\phi(x)dx\\
    = \,\,&\delta,
\end{align*}
where the inequality is derived from $x \geq z_{1-\delta}$. The result follows.

Taking log approximation on both sides, we get
\begin{align*}
    i^L_{t} = i_{t-1} + \left[\delta-\exp\left(\frac{1}{2}\sigma^2_{i,t-1}+z_{1-\delta}\sigma_{i,t-1}\right)\left(1-\Phi\left(z_{1-\delta}+\sigma_{i,t-1}\right)\right)\right].
\end{align*}
It remains to show that the second term on the right-hand side is increasing in $\sigma_{i,t-1}$. Differentiating the second term with respect to $\sigma_{i,t-1}$, we get
\begin{align*}
    &\frac{d}{d\sigma_{i,t-1}}\left[\delta-\exp\left(\frac{1}{2}\sigma^2_{i,t-1}+z_{1-\delta}\sigma_{i,t-1}\right)\left(1-\Phi\left(z_{1-\delta}+\sigma_{i,t-1}\right)\right)\right]\\
     = &\,\exp\left(\frac{1}{2}\sigma^2_{i,t-1}+z_{1-\delta}\sigma_{i,t-1}\right)\left[\phi(z_{1-\delta}+\sigma_{i,t-1}) - (z_{1-\delta}+\sigma_{i,t-1})\left(\Phi(-z_{1-\delta}-\sigma_{i,t-1})\right)\right].
\end{align*}

Consider the function $f(a) = \phi(a) - a\Phi(-a)$. Note that
\begin{align*}
    f'(a) &= \phi'(a) - \Phi(-a) + a\phi(-a)\\
    &= -a\phi(a) -\Phi(-a) + a\phi(a)\\
    &= -\Phi(-a) < 0.
\end{align*}
Moreover, by the Mills ratio,
\begin{align*}
    \lim_{a\to+\infty}
    \left[
    \phi(a)-a\Phi(-a)
    \right]
    =
    0.
\end{align*}
Since $f'(a)=-\Phi(-a)<0$, $f$ is strictly decreasing and converges to zero from above. Therefore $f(a)>0$ for every finite $a$. This implies that the second term on the right-hand side is increasing in $\sigma_{i,t-1}$. Then we are able to define the second term as $s_{i,t} = \bar{s} + \psi \operatorname{Var}_{t-1}\left(R_{i,t}\right)$ in a reduced-form representation of competitive risk-based lending.

Another way is to completely solve for the optimization problem subject to the balance sheet constraint \eqref{eq:balance_sheet_constraint} with Lagrange Multiplier $\lambda^f_t$. This provides the optimal lending amount of the financial intermediary, as is stated in the following Lemma.
\begin{lemma}[Optimal Lending Amount]\label{lemma:optimal_lending_amount}
    Suppose the intermediary chooses $\{\ell_{i,t}\}_{i=1}^N$ to maximize expected profit subject to the balance-sheet constraint and the default-tolerance constraint. Then the optimal lending amount $\ell_{i,t}$ is characterized by
    \begin{align*}
        \Phi(q_{1,i}) =
        \left\{
        \begin{array}{ll}
            \dfrac{1}{1+i^L_t},
            & \text{if } \sum_{i=1}^N \ell_{i,t}<D_{t-1}
            \text{ and the default-tolerance constraint does not bind}, \\[1.2em]
            \dfrac{1+i_{t-1}}{1+i^L_t},
            & \text{if } \sum_{i=1}^N \ell_{i,t}=D_{t-1}
            \text{ and the default-tolerance constraint does not bind}, \\[1.2em]
            1-\delta,
            & \text{if the default-tolerance constraint binds}.
        \end{array}
        \right.
    \end{align*}
\end{lemma}

\begin{proof}
    Let
    \begin{align*}
        m_i(\ell_{i,t})
        :=
        \mathbb{E}_{t-1}
        \left[
        \min\left\{
        R_{i,t},(1+i^L_t)\ell_{i,t}
        \right\}
        \right].
    \end{align*}
    Using the lognormal formula derived above, we can write
    \begin{align*}
        m_i(\ell_{i,t})
        =
        p_{i,t}\mu^A_{i,t|t-1}\tilde{y}_{i,t}
        \left(1-\Phi(q_{2,i})\right)
        +
        (1+i^L_t)\ell_{i,t}\Phi(q_{1,i}).
    \end{align*}
    The derivative of $m_i(\ell_{i,t})$ with respect to $\ell_{i,t}$ is
    \begin{align*}
        m_i'(\ell_{i,t})
        =
        (1+i^L_t)\Phi(q_{1,i}),
    \end{align*}
    where the density terms cancel by Lemma \ref{lemma:pdf_q1_q2}.

    \noindent\textbf{Case 1: Balance-sheet constraint does not bind and the default-tolerance constraint does not bind.}
    When the default-tolerance constraint does not bind, the lending FOC is
    \begin{align*}
        m_i'(\ell_{i,t})
        =
        1+\lambda^f_t,
    \end{align*}
    where $\lambda^f_t$ is the multiplier on the balance-sheet constraint
    \begin{align*}
        D_{t-1}-\sum_{i=1}^N\ell_{i,t}\geq 0.
    \end{align*}
    Therefore,
    \begin{align*}
        (1+i^L_t)\Phi(q_{1,i})
        =
        1+\lambda^f_t.
    \end{align*}
    If the balance-sheet constraint does not bind, then complementary slackness implies
    \begin{align*}
        \lambda^f_t=0.
    \end{align*}
    Hence
    \begin{align*}
        \Phi(q_{1,i})
        =
        \frac{1}{1+i^L_t}.
    \end{align*}

    \noindent\textbf{Case 2: Balance-sheet constraint binds and the default-tolerance constraint does not bind.}
    If the balance-sheet constraint binds, then
    \begin{align*}
        D_{t-1}
        =
        \sum_{i=1}^N \ell_{i,t}.
    \end{align*}
    Therefore, the intermediary's expected profit can be rewritten as
    \begin{align*}
        \mathbb{E}_{t-1}\Pi^f_t
        =
        \sum_{i=1}^N
        \left[
        p_{i,t}\mu^A_{i,t|t-1}\tilde{y}_{i,t}
        \left(1-\Phi(q_{2,i})\right)
        +
        (1+i^L_t)\ell_{i,t}\Phi(q_{1,i})
        -
        (1+i_{t-1})\ell_{i,t}
        \right].
    \end{align*}
    Since the default-tolerance constraint does not bind in this case, the FOC is
    \begin{align*}
        m_i'(\ell_{i,t})
        =
        1+i_{t-1}.
    \end{align*}
    Using $m_i'(\ell_{i,t})=(1+i^L_t)\Phi(q_{1,i})$, we obtain
    \begin{align*}
        (1+i^L_t)\Phi(q_{1,i})
        =
        1+i_{t-1}.
    \end{align*}
    Hence
    \begin{align*}
        \Phi(q_{1,i})
        =
        \frac{1+i_{t-1}}{1+i^L_t}.
    \end{align*}
    Comparing this condition with the original KKT condition
    \begin{align*}
        (1+i^L_t)\Phi(q_{1,i})
        =
        1+\lambda^f_t
    \end{align*}
    also implies that, in this case,
    \begin{align*}
        \lambda^f_t=i_{t-1}.
    \end{align*}

    \noindent\textbf{Case 3: Default-tolerance constraint binds.} The result directly follows from Corollary \ref{corollary:default_tolerance_constraint_binds}.
\end{proof}

This lemma shows that the optimal lending amount is determined by the smaller of two values: the unconstrained lending amount implied by the intermediary's expected-profit FOC and the maximum lending amount allowed by the default-tolerance constraint. When the default-tolerance constraint is slack, the intermediary chooses lending so that the marginal expected repayment equals the relevant marginal funding cost. If the balance-sheet constraint is slack, this marginal cost is one. If the balance-sheet constraint binds, this marginal cost is $1+i_{t-1}$, and the shadow value of balance-sheet capacity in the original KKT formulation is $\lambda^f_t=i_{t-1}$.

When the default-tolerance constraint binds, lending is pinned down by
\begin{align*}
    \Phi(q_{1,i})=1-\delta.
\end{align*}
Equivalently, letting $z_{1-\delta}:=\Phi^{-1}(1-\delta)$, we have
\begin{align*}
    q_{1,i}=z_{1-\delta}.
\end{align*}
Using the definition of $q_{1,i}$, this implies the default-tolerance lending cap
\begin{align*}
    \ell_{i,t}^{\max}
    =
    \frac{
    p_{i,t}\tilde{y}_{i,t}\exp(\bar a_i)
    }
    {1+i^L_t}
    \exp\left(
    -z_{1-\delta}\sigma_{i,t-1}
    \right).
\end{align*}
Equivalently, since
\begin{align*}
    \mu^A_{i,t|t-1}
    =
    \exp\left(
    \bar a_i+\frac{1}{2}\sigma^2_{i,t-1}
    \right),
\end{align*}
we can write
\begin{align*}
    \ell_{i,t}^{\max}
    =
    \frac{
    p_{i,t}\tilde{y}_{i,t}\mu^A_{i,t|t-1}
    }
    {1+i^L_t}
    \exp\left(
    -\frac{1}{2}\sigma^2_{i,t-1}
    -
    z_{1-\delta}\sigma_{i,t-1}
    \right).
\end{align*}
Therefore, when the default-tolerance constraint binds, the optimal lending amount is decreasing in the loan rate $i^L_t$ and in firm-level uncertainty $\sigma_{i,t-1}$. Since the loan rate satisfies $i^L_t=i_{t-1}+s_t$ with $s_t=\bar s+\psi U_t$, higher aggregate uncertainty $U_t$ reduces lending through a higher loan rate in the binding default-tolerance region.
 


\subsubsection{Monetary Authority}

The monetary authority sets the nominal interest rate $i_t$ at the beginning of each period. This rate applies from the beginning of period $t+1$ to its end, equivalently, deposits $D_t$ earn $i_t$ in period $t+1$. Therefore, period-$t$ loans are priced off the predetermined funding rate $i_{t-1}$ through
\begin{align}
    i^L_{i,t}
    =
    i_{t-1}
    +
    s_{i,t}.
\end{align}
In the benchmark model, the nominal interest rate process is exogenous. We introduce an endogenous policy rule when discussing policy implications.




\subsection{Agents' Decisions}
\subsubsection{Households}
Households choose consumption $C_t$, labour supply $N_t$, and deposits $D_t$ to maximized their discounted lifetime utility:
\begin{align*}
    \max_{C_t, N_t, D_t} \mathbb{E}_0\sum_{t=0}^\infty \beta^t U(C_t, N_t),
\end{align*}
where $\beta\in(0,1)$ is the discount factor, subject to the budget constraint \eqref{eq:hh_budget}. The first-order conditions (FOC) imply the Euler equation \eqref{eq:hh_foc_euler} and the labour supply \eqref{eq:hh_foc_labour}:
\begin{align}
    1 &= \beta \mathbb{E}_t \left[\frac{1+i^D_t}{1+\pi_{t+1}}\frac{C_t}{C_{t+1}}\right], \label{eq:hh_foc_euler} \\
    \frac{W_t}{P_t} &= \chi C_t N_t^{\varphi}. \label{eq:hh_foc_labour}
\end{align}

\subsubsection{Firms}
Firms choose labour $n_{i,t}$, intermediate inputs $\{x_{ij,t}\}_{j=1}^N$ and debt $b_{i,t}$ to maximize their nominal profit $\Pi^y_{i,t}$:
\begin{align*}
    \max_{n_{i,t}, \{x_{ij,t}\}_{j=1}^N, b_{i,t}} \mathbb{E}_{t-1}\Pi^y_{i,t},
\end{align*}
subject to the production function \eqref{eq:firm_prod_fcn}, the financial constraint \eqref{eq:firm_financial_constraint} with Lagrange Multiplier $\lambda^1_{i,t}$ and the borrowing constraint \eqref{eq:borrowing_constraint} with Lagrange Multiplier $\lambda^2_{i,t}$. 

Expanding using the definition of $\Pi^y_{i,t}$ and $\tilde{y}_{i,t}$, the objective function can be rewritten as
\begin{align*}
    \mathbb{E}_{t-1}\Pi^y_{i,t} = \mu^A_{i,t|t-1}p_{i,t} \tilde{y}_{i,t}\Phi(q_{2,i}) - (1+i^L_{i,t})b_{i,t} \Phi(q_{1,i}),
\end{align*}
where $\Phi(\cdot)$ is the cumulative distribution function of the standard normal distribution with probability density function $\phi(\cdot)$ and $q_{1,i}$ and $q_{2,i}$ are defined as in the proof of Lemma \ref{lemma:conditional_mean_std_productivity_benchmark}. In equilibrium, financial-market clearing implies $b_{i,t}=\ell_{i,t}$. Therefore the density identity in Lemma \ref{lemma:pdf_q1_q2} applies with $\ell_{i,t}=b_{i,t}$ in the firm problem:
\begin{align*}
    \frac{\phi(q_{2,i})}{\phi(q_{1,i})} = \frac{(1+i^L_{i,t})b_{i,t}}{\mu^A_{i,t|t-1}p_{i,t}\tilde{y}_{i,t}}.
\end{align*}
We can thus simplify the FOCs as follows:
\begin{align}
    [n_{i,t}]:\, &\mu^A_{i,t|t-1}\alpha_{i,t}\frac{\tilde{y}_{i,t}}{n_{i,t}}\Phi(q_{2,i}) = \lambda^1_{i,t}\frac{W_t}{p_{i,t}} - \lambda^2_{i,t}\kappa \mu^A_{i,t|t-1} \exp\left(-\theta \sigma^2_{i,t-1}\right) \alpha_{i,t} \frac{\tilde{y}_{i,t}}{n_{i,t}};\notag \\
    [x_{ij,t}]:\, &\mu^A_{i,t|t-1}\omega_{ij,t}\frac{\tilde{y}_{i,t}}{x_{ij,t}}\Phi(q_{2,i})  = \lambda^1_{i,t}\frac{p_{j,t}}{p_{i,t}} - \lambda^2_{i,t}\kappa \mu^A_{i,t|t-1} \exp\left(-\theta \sigma^2_{i,t-1}\right)  \omega_{ij,t} \frac{\tilde{y}_{i,t}}{x_{ij,t}},\quad \forall j;\notag\\
    [b_{i,t}]:\, &(1+i^L_{i,t})\Phi(q_{1,i}) = \lambda^1_{i,t} - \lambda^2_{i,t}. \label{eq:firm_borrowing_foc}
\end{align}

The first two FOCs implies that the cost-share conditions hold:
\begin{align}
    \frac{W_t n_{i,t}}{p_{j,t}x_{ij,t}} &= \frac{\alpha_{i,t}}{\omega_{ij,t}},\quad \forall j, \label{eq:firm_labour_input} \\
    \frac{p_{j,t}x_{ij,t}}{p_{k,t}x_{ik,t}} &= \frac{\omega_{ij,t}}{\omega_{ik,t}},\quad \forall j \neq k. \label{eq:firm_different_input}
\end{align}
The condition \eqref{eq:firm_borrowing_foc} implies the following Lemma.
\begin{lemma}[Binding Working-Capital Constraint]\label{lemma:working_capital_binding}
    The working-capital constraint \eqref{eq:firm_financial_constraint} binds in any interior solution with finite $q_{1,i}$.
\end{lemma}
\begin{proof}
    From the FOC with respect to $b_{i,t}$,
    \begin{align*}
        \lambda^1_{i,t} - \lambda^2_{i,t} = (1+i^L_{i,t})\Phi(q_{1,i}).
    \end{align*}
    Since $1+i^L_{i,t}>0$ and $\Phi(q_{1,i})>0$ for finite $q_{1,i}$, we have
    \begin{align*}
        \lambda^1_{i,t} = \lambda^2_{i,t} + (1+i^L_{i,t})\Phi(q_{1,i}) > 0.
    \end{align*}
    By complementary slackness for the working-capital constraint,
    \begin{align*}
        \lambda^1_{i,t}\left[b_{i,t}-W_t n_{i,t}-\sum_{j=1}^N p_{j,t}x_{ij,t}\right]=0.
    \end{align*}
    Since $\lambda^1_{i,t}>0$, it follows that
    \begin{align*}
        b_{i,t}=W_t n_{i,t}+\sum_{j=1}^N p_{j,t}x_{ij,t}.
    \end{align*}
\end{proof}



\subsection{Equilibrium}
There are three market in the economy: the goods market, the labour market, and the financial market. The goods market clears when 
\begin{align}\label{eq:goods_market_clear}
    y_{i,t} = c_{i,t} + \sum_{j=1}^N x_{ji, t}, \; \forall i.
\end{align}
This means that the total output of firm $i$ is equal to the final comsumption of the household plus the intermediate inputs used by other firms.

The labour market clears when
\begin{align}\label{eq:labour_market_clear}
    N_t = \sum_{i=1}^N n_{i,t}.
\end{align}

The financial market clears when
\begin{align}\label{eq:financial_market_clear}
    b_{i,t} = \ell_{i,t}, \forall i.
\end{align}

Given initial nominal deposits $D_{-1}$, initial money supply $M^S_{-1}$,
and exogenous policy processes $\{i_t,X_t\}_{t=0}^{\infty}$, a competitive equilibrium is a stochastic sequence of allocations
\begin{align*}
    \left\{
    C_t,\{c_{i,t}\}_{i=1}^N,N_t,D_t,
    \{n_{i,t}\}_{i=1}^N,
    \{x_{ij,t}\}_{i,j=1}^N,
    \{y_{i,t}\}_{i=1}^N,
    \{b_{i,t}\}_{i=1}^N,
    \{\ell_{i,t}\}_{i=1}^N
    \right\}_{t=0}^{\infty},
\end{align*}
prices and rates
\begin{align*}
    \left\{
    P_t,\{p_{i,t}\}_{i=1}^N,W_t,
    i^D_t,\{i^L_{i,t}\}_{i=1}^N,
    \{s_{i,t}\}_{i=1}^N
    \right\}_{t=0}^{\infty},
\end{align*}
and profits and policy variables
\begin{align*}
    \left\{
    \Pi^y_t,\Pi^f_t,M^S_t,X_t
    \right\}_{t=0}^{\infty}
\end{align*}
such that, for all $t$ and all states:

\begin{enumerate}
    \item Households maximize expected lifetime utility subject to the nominal budget constraint.
    \item Firms maximize expected nominal profits subject to the production function, the working-capital constraint, and the borrowing constraint.
    \item Financial intermediaries satisfy the balance-sheet constraint, the loan-pricing rule, and the borrowing-capacity rule.
    \item The monetary authority sets $i_t$ according to the benchmark policy process.
    \item Goods, labor, and financial markets clear.
\end{enumerate}

\subsection{An Economy without Production Network}
We would like to compare the effect of uncertainty shock between a network economy with a non-networked economy. We assume the same household with the same objective and budget constraint. We also adapt the same environment for financial intermediary and monetary authority. However, the firms have a different setting.

To make the production process comparable, we assume that the production of a variety of good $i$ requires some capital $k_{i,t}$ instead of the intermediate goods with the following production function
\begin{align*}
    y_{i,t} = A_{i,t} n_{i,t}^{\alpha_{i,t}}k_{i,t}^{1-\alpha_{i,t}}.
\end{align*}
The capital is assumed to be borrowed at the beginning of each period from the financial intermediary with price $p^k_{i,t}$. Therefore, the budget constraint of the firms is
\begin{align*}
    W_t n_{i,t} + p^k_{i,t} k_{i,t} \leq b_{i,t}.
\end{align*}

At the same time, the goods market clears when 
\begin{align*}
    y_{i,t} = c_{i,t}.
\end{align*}








\clearpage
\bibliographystyle{aer}
\bibliography{network_uncertainty}


\clearpage


%%%%%%%%%%%%%%%%%%%%%%%%%%
% Publication Appendix
%%%%%%%%%%%%%%%%%%%%%%%%%%
%%%%%%%%%%%%%%%%%%%%%%%%%%%%%%%%%%%%%%%%%%%%%%%%%%%%%%%%%%%%%%%%%%%%%%%%%%%%%%%%%%%%%%%
%%%%%%%%%%%%%%%%%%%%%%%%%%%%%%%%%%%%%%%%%%%%%%%%%%%%%%%%%%%%%%%%%%%%%%%%%%%%%%%%%%%%%%%
%%%%%%%%%%%%%%%%%%%%%%%%%%%%%%%%%%%%%%%%%%%%%%%%%%%%%%%%%%%%%%%%%%%%%%%%%%%%%%%%%%%%%%%
%%%%%%%%%%%%%%%%%%%%%%%%%%%%%%%%%%%%%%%%%%%%%%%%%%%%%%%%%%%%%%%%%%%%%%%%%%%%%%%%%%%%%%%
% \clearpage
\begin{appendix}
% Formatting options
\addcontentsline{toc}{section}{Appendices}
\renewcommand{\thesection}{\Roman{section}} 
\renewcommand{\thesubsection}{\Roman{subsection}}
\renewcommand{\theequation}{\thesection\arabic{equation}}
\numberwithin{equation}{section}
% \setcounter{page}{1}
\setcounter{figure}{0} \renewcommand{\thefigure}{\Roman{section}.\arabic{figure}}
\setcounter{table}{0} \renewcommand{\thetable}{\Roman{section}.\arabic{table}}
\setcounter{subsection}{0} \renewcommand{\thesubsection}{\Roman{section}.\arabic{subsection}}
\setcounter{equation}{0} \renewcommand{\theequation}{\Roman{section}.\arabic{equation}}
\setcounter{theorem}{0} \renewcommand{\thetheorem}{\Roman{section}.\arabic{theorem}}
\setcounter{definition}{0} \renewcommand{\thedefinition}{\Roman{section}.\arabic{definition}}
% \setcounter{proposition}{0} \renewcommand{\theproposition}{\Roman{section}.\arabic{proposition}}
\setcounter{assumption}{0} \renewcommand{\theassumption}{\Roman{section}.\arabic{assumption}}
% \setcounter{hyp}{0} \renewcommand{\thehyp}{\Roman{section}.\arabic{hyp}}
\setcounter{footnote}{0}
%%%%%%%%%%%%%%%%%%%%%%%%%%%%%%%%%%%%%%%%%%%%%%%%%%%%%%%%%%%%%%%%%%%%%%%%%%%%%%%%%%%%%%%
%%%%%%%%%%%%%%%%%%%%%%%%%%%%%%%%%%%%%%%%%%%%%%%%%%%%%%%%%%%%%%%%%%%%%%%%%%%%%%%%%%%%%%%
%%%%%%%%%%%%%%%%%%%%%%%%%%%%%%%%%%%%%%%%%%%%%%%%%%%%%%%%%%%%%%%%%%%%%%%%%%%%%%%%%%%%%%%
%%%%%%%%%%%%%%%%%%%%%%%%%%%%%%%%%%%%%%%%%%%%%%%%%%%%%%%%%%%%%%%%%%%%%%%%%%%%%%%%%%%%%%%



\end{appendix}
\end{document}